\documentclass[12pt, a4paper]{article}
\usepackage[english]{babel}
\usepackage[utf8]{inputenc}
\usepackage[T1]{fontenc}
\usepackage{amsmath, amssymb, amsthm}
\usepackage{geometry}
\usepackage{hyperref}
\usepackage{listings}

\geometry{margin=1in}

\newtheorem{theorem}{Theorem}
\newtheorem{lemma}{Lemma}
\newtheorem{definition}{Definition}

\title{On the Erd\H{o}s-Hajnal Conjecture: Structure and Independence in $H$-Free Graphs}
\author{Charles EDOU NZE\thanks{Charles EDOU NZE, chercheur indépendant}}
\date{\today}

\begin{document}
\maketitle

\begin{abstract}
We present a detailed, zero-ellipse exposition of the Erd\H{o}s-Hajnal conjecture and its foundational consequences. This document formally defines the problem, constructs an architecture for autoformalization in systems such as Lean 4, and proceeds with a rigorous step-by-step mathematical proof demonstrating bounds on homogeneous sets in graphs avoiding specific induced subgraphs.
\end{abstract}

\tableofcontents
\newpage

\section{Axiomatic Definitions and Problem Formulation}

Let $G = (V, E)$ be a finite, simple, undirected graph.

\begin{definition}
A graph $G$ contains a graph $H$ as an induced subgraph if there exists a subset of vertices $S \subseteq V(G)$ such that the subgraph of $G$ induced by $S$, denoted $G[S]$, is isomorphic to $H$. If $G$ does not contain $H$ as an induced subgraph, $G$ is called $H$-free.
\end{definition}

\begin{definition}
A clique in $G$ is a subset of vertices $C \subseteq V$ such that for all $u, v \in C$ with $u \neq v$, the edge $\{u, v\}$ is in $E$. An independent set is a subset $I \subseteq V$ such that for all $u, v \in I$, $\{u, v\} \notin E$. A homogeneous set is either a clique or an independent set. The size of the largest homogeneous set in $G$ is denoted $h(G) = \max(\omega(G), \alpha(G))$, where $\omega(G)$ is the clique number and $\alpha(G)$ is the independence number.
\end{definition}

The Erd\H{o}s-Hajnal Conjecture states:

\begin{theorem}[Erd\H{o}s-Hajnal]
For every finite graph $H$, there exists a constant $c(H) > 0$ such that for every $H$-free graph $G$ on $n$ vertices,
\begin{equation}
h(G) \geq n^{c(H)}
\end{equation}
\end{theorem}

\section{Architecture for Autoformalization (Lean 4)}

\begin{lstlisting}[language=Python, basicstyle=\ttfamily\small]
import Mathlib.Combinatorics.SimpleGraph.Basic
import Mathlib.Combinatorics.SimpleGraph.Clique

universe u v

variables {V : Type u} [Fintype V]
variables (G : SimpleGraph V) (H : SimpleGraph V)

def IsInducedSubgraph (H G : SimpleGraph V) : Prop :=
  \exists (S : Set V), Nonempty (G.induce S \simeqg H)

def HFree (H G : SimpleGraph V) : Prop :=
  \not IsInducedSubgraph H G

def HomogeneousSet (G : SimpleGraph V) (S : Set V) : Prop :=
  G.IsClique S \lor (G.compl).IsClique S

theorem Erdos_Hajnal_Conjecture (H : SimpleGraph V) :
  \exists (c : \mathbb{R}), c > 0 \land
  \forall (G : SimpleGraph V), HFree H G \rightarrow
  \exists (S : Set V), HomogeneousSet G S \land
  (Fintype.card S : \mathbb{R}) \geq (Fintype.card V : \mathbb{R})^c
\end{lstlisting}

\section{Contextual Literature Research}

The problem is fundamentally linked to Ramsey theory. Ramsey's theorem states that any graph on $n$ vertices contains a homogeneous set of size at least $\frac{1}{2} \log_2 n$. The Erd\H{o}s-Hajnal conjecture postulates that excluding any induced subgraph $H$ drastically changes this behavior, forcing a polynomial-sized homogeneous set.

A profound analogy exists with the perfect graph theorem and the strong perfect graph theorem (Chudnovsky, Robertson, Seymour, Thomas, 2006). Perfect graphs (which exclude odd holes and odd antiholes) satisfy $h(G) \geq n^{1/2}$. The conjecture asks whether excluding just one arbitrary finite graph $H$ is sufficient to bound $h(G)$ polynomially.

\section{Proof Strategy and Lemmata}

\subsection{Lemma 1: The Base Case for Small $H$}

\begin{lemma}
If $H$ is a graph on 2 vertices, the conjecture holds with $c(H) = 1$.
\end{lemma}

\begin{proof}
Let $H$ be a graph on 2 vertices. There are two such graphs: $K_2$ (a single edge) and $\overline{K_2}$ (two non-adjacent vertices).

Case 1: $H = K_2$.
Suppose $G$ is a $K_2$-free graph on $n$ vertices. By definition, $G$ contains no edges. Therefore, the entire vertex set $V(G)$ is an independent set. Thus, $\alpha(G) = n$.
We have $h(G) = \max(\omega(G), \alpha(G)) = \max(1, n) = n = n^1$.
In this case, $c(K_2) = 1 > 0$.

Case 2: $H = \overline{K_2}$.
Suppose $G$ is an $\overline{K_2}$-free graph on $n$ vertices. By definition, $G$ contains no two non-adjacent vertices. Therefore, every pair of distinct vertices is connected by an edge. This means $G$ is a complete graph $K_n$.
The entire vertex set $V(G)$ is a clique. Thus, $\omega(G) = n$.
We have $h(G) = \max(\omega(G), \alpha(G)) = \max(n, 1) = n = n^1$.
In this case, $c(\overline{K_2}) = 1 > 0$.

In both cases, for any graph $H$ on $|V(H)| = 2$, there exists $c(H) = 1 > 0$ such that $h(G) \geq n^{c(H)}$.
\end{proof}

\subsection{Lemma 2: Ramsey's Bound is Insufficient}

\begin{lemma}
Without the $H$-free condition, there exist graphs for which $h(G) = O(\log n)$.
\end{lemma}

\begin{proof}
We apply the probabilistic method of Erd\H{o}s. Consider a random graph $G \in \mathcal{G}(n, 1/2)$, where each edge is included with probability $p = 1/2$ independently.

Let $k$ be an integer. The probability that a specific subset of $k$ vertices forms a clique is $(1/2)^{\binom{k}{2}}$.
Similarly, the probability that a specific subset of $k$ vertices forms an independent set is $(1/2)^{\binom{k}{2}}$.

Let $X$ be the random variable representing the number of homogeneous sets of size $k$. By linearity of expectation:
\begin{align*}
\mathbb{E}[X] &= \binom{n}{k} \left( \left(\frac{1}{2}\right)^{\binom{k}{2}} + \left(\frac{1}{2}\right)^{\binom{k}{2}} \right) \\
&= 2 \binom{n}{k} 2^{-\binom{k}{2}} \\
&\leq 2 \frac{n^k}{k!} 2^{-k(k-1)/2}
\end{align*}

If we set $k = 2 \log_2 n + 1$, then:
\begin{align*}
2^{-k(k-1)/2} &= 2^{-\frac{(2 \log_2 n + 1)(2 \log_2 n)}{2}} \\
&= 2^{-2(\log_2 n)^2 - \log_2 n} \\
&= n^{-2\log_2 n - 1}
\end{align*}

Thus,
\begin{align*}
\mathbb{E}[X] &\leq 2 \frac{n^{2\log_2 n + 1}}{k!} n^{-2\log_2 n - 1} \\
&= \frac{2}{k!} < 1
\end{align*}

Since the expected number of homogeneous sets of size $k = 2\log_2 n + 1$ is strictly less than 1, there must exist at least one graph in the probability space with exactly 0 homogeneous sets of size $k$.
For this graph, $h(G) < 2 \log_2 n + 1$, meaning $h(G) = O(\log n)$.

This demonstrates that the assumption of being $H$-free is strictly necessary to force a polynomial lower bound $n^{c(H)}$.
\end{proof}

\section{Conclusion}

The base cases and probabilistic bounds establish the fundamental tension of the Erd\H{o}s-Hajnal conjecture. While general graphs contain only logarithmic homogeneous sets, the exclusion of any finite structure $H$ is theorized to force a massive macroscopic structural regularity, expanding the homogeneous set size exponentially relative to the generic case. The isolation of these properties sets the stage for further decomposition via the substitution method and modular decomposition.

\end{document}
